\section{Lecture 17 (February 24, 2025, est.): The Verlinde category}
\label{lec17}
\noindent\textit{Scribe: Guillermo Sanmarco}

Throughout this lecture, $k$ is algebraically closed of characteristic
$p>0$. We apply semisimplification to representations of the cyclic group
$C_p$. Recall that an indecomposable survives semisimplification precisely
when its categorical dimension is nonzero, provided traces of nilpotent
endomorphisms vanish. This hypothesis holds in representation categories,
where the categorical trace is the ordinary linear trace.

\subsection{Representations of the cyclic group}

Choose a generator $\sigma$ of $C_p$. Writing $t=\sigma-1$ gives
\[
 k[C_p]=k[\sigma]/(\sigma^p-1)\cong k[t]/(t^p).
\]
Thus a representation is a vector space with a nilpotent endomorphism
whose $p$th power is zero. Jordan normal form gives exactly $p$
indecomposable representations:
\[
 J_n=k[t]/(t^n),\qquad 1\leq n\leq p.
\]
In particular, $J_1$ is the trivial representation and $J_p$ is the
regular representation.

\begin{definition}
The \emph{Verlinde category} is the symmetric tensor category
\[
 \operatorname{Ver}_p=\overline{\Rep_k(C_p)}.
\]
Write $L_n$ for the image of $J_n$ in this semisimplification. Its simple
objects are $L_1,\ldots,L_{p-1}$, with $L_1=\1$, whereas $L_p=0$.
Their categorical dimensions are $\dim(L_n)=n\cdot1_k$.
\end{definition}

There are finitely many simples here, so $\operatorname{Ver}_p$ is a
symmetric fusion category. Finiteness of the original group alone would
not ensure this conclusion for its semisimplification: representation
categories of other finite groups can have infinitely many surviving
indecomposables. See \cite{EOSemisimplification} for examples and the
general construction.

\subsection{Tensor products of Jordan blocks}

\begin{proposition}\label{prop:lec17-jordan-products}
The $J_n$ are self-dual, and
\[
 J_2\otimes J_n\cong
 \begin{cases}
 J_2,&n=1,\\
 J_{n-1}\oplus J_{n+1},&2\leq n\leq p-1,\\
 J_p^{\oplus2},&n=p.
 \end{cases}
\]
For $1\leq n\leq p-1$ one also has
\[
 J_{p-1}\otimes J_n\cong J_{p-n}\oplus J_p^{\oplus(n-1)}.
\]
\end{proposition}
\begin{proof}
On a dual representation, $\sigma$ acts by the inverse transpose.
The inverse of a unipotent Jordan block is again a unipotent matrix with
one Jordan block of the same size, proving self-duality.

For the first tensor-product formula, identify the underlying vector
space with $k[x,y]/(x^2,y^n)$. The generator acts by multiplication by
$(1+x)(1+y)$. Set $u=x(1+y)$; this is an invertible change of variable,
so the algebra is also $k[u,y]/(u^2,y^n)$, and $t$ acts by multiplication
by $u+y$. Hence
\[
 t^r=y^r+r u y^{r-1}.
\]
For $2\leq n<p$, the nilpotency index is $n+1$, since
$t^n=nuy^{n-1}\ne0$. Moreover, the kernel of multiplication by $t$
has dimension two. Indeed, $t(A(y)+uB(y))=0$ means
\[
 yA=0,\qquad A+yB=0,
\]
so $A=ay^{n-1}$ and $B=-ay^{n-2}+by^{n-1}$. There are therefore two
Jordan blocks; their sizes are $n+1$ and $n-1$ because the total
dimension is $2n$. If $n=p$, the same kernel computation gives two
blocks, each of size at most $p$, and their total size is $2p$.
The case $n=1$ is immediate.

Finally, tensor the exact sequence
\[
 0\longrightarrow J_{p-1}\longrightarrow J_p
 \longrightarrow J_1\longrightarrow0
\]
with $J_n$. The middle term is $J_p^{\oplus n}$: for any finite group,
the regular representation tensored diagonally with $V$ is isomorphic
to $\dim(V)$ copies of the regular representation, by the map
$g\otimes v\mapsto g\otimes g^{-1}v$.
A surjection $(k[t]/t^p)^n\to k[t]/t^n$ can be changed by an invertible
basis transformation of its source to the standard quotient on its
first factor and zero on the others. To see this, one of the images of
the standard generators must have nonzero constant term; make it the
first generator, multiply it by a unit to make its image $1$, and
subtract suitable multiples of it from the remaining generators.
The kernel is then
$t^n(k[t]/t^p)\oplus(k[t]/t^p)^{n-1}$, as claimed.
\end{proof}

\begin{theorem}[Fusion rules]\label{thm:lec17-fusion}
For $1\leq i,j\leq p-1$,
\[
 L_i\otimes L_j\cong
 \bigoplus_{r=1}^{\min(i,j,p-i,p-j)}L_{|i-j|+2r-1}.
\]
In particular,
\[
 L_2\otimes L_n=L_{n-1}\oplus L_{n+1},\qquad
 L_{p-1}\otimes L_n=L_{p-n},
\]
where $L_0=L_p=0$ in the first formula.
\end{theorem}
\begin{proof}
The two displayed special cases follow from the proposition. For
$p>2$, write $x=[L_2]$ in the Grothendieck ring and define polynomials
$f_0=0$, $f_1=1$, $f_{n+1}=xf_n-f_{n-1}$. The recurrence gives
$[L_n]=f_n(x)$ and $f_p(x)=0$.
Induction using the same recurrence gives the polynomial identity
\[
 f_i(x)f_j(x)=\sum_{r=1}^{\min(i,j)}f_{|i-j|+2r-1}(x).
\]
Modulo $f_p(x)$ one has $f_{p+a}(x)=-f_{p-a}(x)$ for $a\geq0$ in
the range in question. The terms above $p$ consequently cancel the
reflected terms below $p$, and the term $f_p$ vanishes. The remaining
indices start at $|i-j|+1$ and end at
$\min(i+j-1,2p-i-j-1)$, giving the formula. Since the category is
semisimple, equality of these classes is an isomorphism of objects.
For $p=2$ there is only the simple unit, and the formula is immediate.
\end{proof}

These are the truncated Clebsch--Gordan rules also occurring in the
semisimplified quantum $\mathfrak{sl}_2$ category at a primitive
$2p$th root of unity in characteristic zero. This comparison concerns
the fusion rules; the categories have different ground-field
characteristics. For the positive-characteristic construction and
fusion rules, see \cite[Section 3]{OstrikVerlinde}.

\subsection{Small characteristics and failure of fiber functors}

For $p=2$, the only surviving simple is the unit, so
$\operatorname{Ver}_2\simeq\Vect_k$. For $p=3$, the simples are $L_1$
and $L_2$, with $L_2^{\otimes2}=\1$. Its self-symmetry is $-1$:
in $J_2\otimes J_2$, the alternating line is a copy of $J_1$ on which
the flip is $-1$, whereas the symmetric square is $J_3$ and disappears.
Thus $\operatorname{Ver}_3\simeq\SVec_k$.

For $p=5$, the fusion rule gives
\[
 L_3\otimes L_3\cong\1\oplus L_3.
\]
If there were a symmetric fiber functor to $\SVec_k$, the total vector
space dimension $n$ of the image of $L_3$ would be a positive integer
satisfying $n^2=n+1$, which is impossible. Here one must use an ordinary
nonnegative integer dimension, not the categorical dimension in $k$:
the categorical dimension $3\in k$ does satisfy this equation in
characteristic $5$.

\begin{proposition}
For every prime $p\geq5$, $\operatorname{Ver}_p$ has no fiber functor
to $\Vect_k$ or to $\SVec_k$.
\end{proposition}
\begin{proof}
Suppose such a functor exists, and let $d_i$ be the total dimension of
the image of $L_i$. Faithfulness gives positive integers $d_i$ for
$1\leq i\leq p-1$. Set $d_0=d_p=0$, $d_1=1$, and $d=d_2$.
The fusion rules imply
\[
 d d_i=d_{i-1}+d_{i+1}\quad(1\leq i\leq p-1).
\]
Since $d^2=1+d_3\geq2$, one has $d\geq2$. Starting with
$d_1-d_0=1$, the recurrence gives
\[
 d_{i+1}-d_i=(d-2)d_i+(d_i-d_{i-1})\geq1
\]
by induction, including the last step $i=p-1$. Thus $d_p\geq p$,
contrary to $d_p=0$.
\end{proof}

\begin{example}[A characteristic-two obstruction]
There are also symmetric tensor categories in characteristic $2$ with
no fiber functor to $\Vect_k$, even though $\operatorname{Ver}_2$ itself
is $\Vect_k$. Consider the Hopf algebra
\[
 H=k[d]/(d^2),\qquad \Delta(d)=d\otimes1+1\otimes d,
\]
with triangular structure $R=1\otimes1+d\otimes d$. The identities
$(\Delta\otimes\id)(R)=R_{13}R_{23}$ and
$(\id\otimes\Delta)(R)=R_{13}R_{12}$ follow by expanding and using
$d^2=0$; also $R_{21}R=1$ because $2=0$. Thus the symmetry on
$\Rep(H)$ is $c=\tau\circ R$, where $\tau$ is the usual flip.

Let $X$ have basis $e,f$ with $de=0$ and $df=e$. Then
\[
 c(e\otimes e)=e\otimes e,\qquad
 c(e\otimes f)=f\otimes e,\qquad
 c(f\otimes f)=f\otimes f+e\otimes e.
\]
The image of $\id-c$ is spanned by $e\otimes e$ and
$e\otimes f+f\otimes e$. Consequently
$\Sym^2X=\operatorname{coker}(\id-c)$ is two-dimensional, with basis
the classes of $e\otimes f$ and $f\otimes f$. The action of $d$ on
both classes is zero, so $\Sym^2X\cong\1^{\oplus2}$.
On the other hand, $X$ is an extension of $\1$ by $\1$. An exact
symmetric fiber functor $F$ to $\Vect_k$ would give $\dim F(X)=2$
and preserve this cokernel. This would imply both
$\dim F(\Sym^2X)=2$ and
$\dim\Sym^2(F(X))=3$, a contradiction.
\end{example}

\subsection{Tensor subcategories of the Verlinde category}

Assume now that $p\geq5$. Let $\operatorname{Ver}_p^+$ be the full
semisimple subcategory with simples
\[
 L_1,L_3,L_5,\ldots,L_{p-2}.
\]
The fusion rules show that it is closed under tensor products and
duals, and that it is generated by $L_3$. More explicitly,
$L_3\otimes L_{2r+1}$ contains $L_{2r+3}$ whenever
$3\leq2r+3\leq p-2$, starting with $L_3$ itself.

The object $\varepsilon=L_{p-1}$ is invertible of order two, with
$\dim(\varepsilon)=-1$. For an invertible object in a symmetric
category, the self-symmetry is its categorical dimension times the
identity (compute the trace using its inverse object). Thus the
subcategory generated by $\varepsilon$ is $\SVec_k$, with
$\varepsilon$ corresponding to the odd line.

\begin{proposition}\label{prop:lec17-subcategories}
The full tensor subcategories of $\operatorname{Ver}_p$ are exactly
\[
 \Vect_k,\qquad \SVec_k,\qquad
 \operatorname{Ver}_p^+,\qquad\operatorname{Ver}_p.
\]
Moreover, multiplication induces a symmetric equivalence
\[
 \SVec_k\boxtimes\operatorname{Ver}_p^+
 \xrightarrow{\ \sim\ }\operatorname{Ver}_p.
\]
Here a full tensor subcategory is understood to be closed under finite
direct sums and subobjects, and $\boxtimes$ denotes the product whose
simple objects are pairs of simple objects of its factors.
\end{proposition}
\begin{proof}
If a tensor subcategory contains $L_i$ with $2\leq i\leq p-2$,
then $L_i\otimes L_i$ contains $L_3$, so the subcategory contains
$\operatorname{Ver}_p^+$. If $i$ is even, then $p-i$ is odd, hence
$L_{p-i}$ already belongs to $\operatorname{Ver}_p^+$. The product
$L_i\otimes L_{p-i}$ contains $L_{p-1}=\varepsilon$, so the
subcategory is all of $\operatorname{Ver}_p$. More generally,
adjoining $\varepsilon$ to $\operatorname{Ver}_p^+$ gives all the
even-indexed simples by $\varepsilon\otimes L_j=L_{p-j}$ for odd $j$.
If there are no such
intermediate simples, the only possibilities are $\Vect_k$ and the
subcategory generated by $\varepsilon$. This proves the list.

The multiplication functor in the displayed equivalence sends
$\1\boxtimes L_j$ to $L_j$ and
$\varepsilon\boxtimes L_j$ to $L_{p-j}$ for odd $j$. These are all
the simple objects of $\operatorname{Ver}_p$, each appearing once.
It identifies the corresponding morphism spaces, and is therefore
an equivalence of semisimple categories. The symmetry allows the
middle two factors of a tensor product to be interchanged, supplying
its symmetric monoidal structure.
\end{proof}

This decomposition and the list of subcategories are also proved in
\cite[Proposition 3.3]{OstrikVerlinde}. The product construction and its
universal property will be made explicit in the next lecture.
